\chapter{The Cancellation}
% OUTLINE Ch3 "The Cancellation (algebra's face: the identity structure)".
% Laws: INTUITION LAW (labeled blocks, adds-no-claim), CITATION (the
% schoolboy Gauss pairing story w/ "by the account" honesty; Black 1946
% mutilated chessboard; Noether 1918 as the summit, her story at
% Cavendish grade; the Ch1 bridge per Kodo's ruling -- the theorem where
% the conservation lands). The series' null method referenced in ONE
% clause (method, not apparatus -- device dose stays Ch2+Ch4). The
% exactness FENCE (a cancellation proves shared content, never a
% self-auditing world). Gate: Kodo.

\begin{intuition}
A bookkeeper closes the month with two ledgers kept two ways --- one by
account, one by day --- and subtracts. The zero that comes back is not a
relief; it is information. Scattered errors do not conspire to cancel to
the coin: an exact zero says the two tellings were tellings of one till.
The feeling to carry: an exact cancellation is never an accident, and a
proof can be built out of exactly that never.
\end{intuition}

Algebra's face of proof is the oldest and the most familiar, so familiar
the reader may never have noticed it is a structure at all: the identity.
Two expressions, transformed step by lawful step into one another; a sum
whose neighbors annihilate in pairs until only the ends remain --- the
schoolroom's telescoping sum; the equality that holds not approximately,
not at last, but exactly and at every instant. The legend of the
structure is a schoolboy: Gauss, set with his class to sum the numbers
from one to a hundred, and answering --- by the account that has come
down from his biographer --- almost at once, having seen that the sum
pairs with itself: first with last, second with second-to-last, fifty
pairs of a hundred and one. The story's lesson is the face's lesson. The
child did not compute harder; he found the \emph{arrangement} in which
the answer is visible --- and an identity is precisely that: a
rearrangement under which content shows itself unchanged.

The structure's power turns destructive --- in the honest sense --- when
the exactness is used to prove impossibility. Take the standard specimen,
published as a puzzle by Max Black in 1946: a chessboard with two
opposite corners removed, and a supply of dominoes each covering two
adjacent squares. Can the mutilated board be tiled? Count colors: every
domino, wherever it lies, covers one dark square and one light; the two
removed corners were the same color; the remaining board has thirty-two
of one shade and thirty of the other. No arrangement of dominoes ---
none, ever, however ingenious --- can cover unequal numbers with pieces
that pair them. The proof visits no arrangements. It exhibits a quantity
that every allowed placement preserves exactly --- the color balance ---
and observes that the goal state disagrees with the start. Where the
last chapter's invariant forced existence, this face's exactness forces
\emph{impossibility}: the ledger cannot be made to balance, because
every legal entry is balanced and the books were opened unbalanced.

\begin{intuition}
The feeling for it: pairing as seeing. The child Gauss did not add; he
folded the sum in half and watched the terms marry. Every proof on this
face has that flavor --- find the pairing, the coloring, the two
tellings, and let the exactness do the work the labor would have done.
\end{intuition}

The face's summit is a theorem that joins this chapter to the first, and
it was proved by the person Hilbert had to fight his own university to
hire. Emmy Noether came to G\"ottingen in 1915, invited by Hilbert and
Klein to untangle a puzzle at the heart of the new general relativity
--- the status of energy conservation --- while the faculty senate
refused her the right to lecture on the grounds of her sex, and Hilbert
made the reply that has outlived every man who voted: they were a
university, not a bathhouse. Her habilitation came in 1919; the theorem
this chapter needs came the year before, in the 1918 paper
\emph{Invariante Variationsprobleme}, and it says: every continuous
symmetry of a variational functional forces a conserved quantity. If the
functional of Chapter 1 is unchanged by shifts of time, a quantity ---
the one the trade calls energy --- is exactly conserved along every
extremal; unchanged by shifts of place, momentum; by rotations, angular
momentum. The reader should see both faces meet in it: the symmetry is a
variational fact --- the first variation of the functional along the
symmetry's direction vanishes identically --- and the conservation is
the exact cancellation that fact forces in the bookkeeping of every
motion. The zero of energy accounting is not a fortunate habit of
nature's ledgers; it is an identity, with a date and an author, riding
on a symmetry. Conservation is cancellation with a reason, and the
reason is Chapter 1's.

And so the fence, one paragraph as the law requires. What an exact
cancellation proves is that two tellings shared content --- that the
functional really had the symmetry, that the coloring really was
preserved, that the two ledgers really described one till. It does not
prove that the world audits itself, prefers balance, or keeps books; the
bookkeeper is the prover, and the zero certifies the \emph{descriptions'}
agreement, always. The reader of this series has watched this face work
an honest shift already --- the first volume's null method, two spellings
of one content subtracted to a zero with the remainder read as signal,
is this chapter's structure practiced as an instrument --- and that one
clause is all this chapter will say of it. Exactness belongs to the
algebra. What the algebra proves is as good as its tellings, and the
label on the feeling --- \emph{the zero is information} --- must never
harden into \emph{the universe balances its accounts}.

What the cancellation cannot do is reach an answer that no finite
rearrangement lands. Some quantities are not equal to any telling's
closed form --- they can only be approached, cornered, pinned between
tellings that agree ever more closely and never coincide. Proving things
about what you can corner but not catch is the fourth face, the
numerical one, and the reader of this series will recognize its gait
before the next chapter names it: it is the face the whole first volume
walked on.
